%% sections/related_work.tex

\section{Related Work}
\label{sec:relatedwork}

\subsection{Zero-Shot and Perplexity-Based Detection}

A foundational observation underlying machine-text detection is that language models assign higher probability to their own outputs than to human text, resulting in characteristically low perplexity scores. DetectGPT~\cite{mitchell:2023} exploits this by examining the local curvature of the log-probability surface: model-generated text tends to occupy local maxima, whereas human text lies in lower-probability regions. Fast-DetectGPT~\cite{bao:2023} improves on this by using conditional probability curvature, reducing computational cost without sacrificing accuracy. Binoculars~\cite{hans:2024} takes a different approach, measuring the ratio between per-token perplexity and cross-perplexity computed by two related models; it achieves near-perfect detection on clean text but degrades under adversarial paraphrasing. DetectLLM~\cite{su:2023} similarly leverages log-rank information. These zero-shot methods are appealing because they require no labeled training data, but they rely on LM probability estimates that can be disrupted by even simple text transformations~\cite{sadasivan:2023}.

\subsection{Supervised Discriminative Classifiers}

Supervised detectors fine-tune pre-trained language models on labeled human/AI pairs. Early work showed that even classifiers trained on GPT-2 outputs generalize well to other generators~\cite{solaiman:2019,ippolito:2020}. More recent detectors based on RoBERTa~\cite{liu:2019} and BERT~\cite{devlin:2019} demonstrate strong in-distribution accuracy, but exhibit systematic degradation on out-of-distribution generators and adversarial inputs. Bevendorff et al.~\cite{bevendorff:2025b} provide a careful analysis of the two main detection paradigms—authorship attribution (assigning text to a known model) and authorship verification (deciding whether a pair of texts share an author)—and argue that the verification framing is better suited to real-world conditions where the target model is unknown. Within PAN, the successive editions of the Voight-Kampff challenge~\cite{bevendorff:2024d,bevendorff:2025d} have documented the rapid escalation of adversarial techniques and the corresponding need for more robust classifiers.

\subsection{Stylometric Approaches}

Stylometry—the quantitative analysis of writing style—was originally developed for human authorship attribution, where it has a long and successful history. The question of whether stylometric signals transfer to AI detection is more nuanced. Schuster et al.~\cite{schuster:2020} demonstrated that classical stylometric features such as function-word distributions and part-of-speech n-grams are insufficient on their own to detect machine-generated fake news, primarily because modern LLMs are capable of approximating human lexical patterns. Nevertheless, several studies have found that certain features—especially those capturing discourse structure and long-range coherence—differ systematically between human and AI text~\cite{uchendu:2020}. The PAN task itself uses a TF-IDF + linear SVM baseline as a reference, which can be seen as a form of bag-of-words stylometry.

Our work combines stylometric analysis with deep semantic representations, an approach that has not been systematically explored in the context of the PAN Voight-Kampff task. The 38 handcrafted features in SLPE draw on established measures of vocabulary richness (TTR, Honore's R~\cite{honore:1979}, MATTR), syntactic complexity, discourse coherence, and word-level entropy, and cover five complementary dimensions of writing style. Crucially, we compute these features entirely without external language models, making the extractor fast, lightweight, and adversarially transparent.

\subsection{Multimodal and Fusion Architectures}

Fusion of heterogeneous feature types has been widely studied in related NLP tasks. Cross-attention mechanisms~\cite{vaswani:2017} have proven effective for late fusion when the two modalities operate at different granularities. In RADAR we apply a single-step cross-attention between the \texttt{[CLS]} representation of ModernBERT and the projected stylometric vector, allowing the semantic stream to selectively attend to the stylometric evidence before the two are combined. This design is deliberately lightweight: the fusion adds fewer than five million trainable parameters on top of the 395-million parameter encoder, making the first training stage computationally efficient while still benefiting from full pre-trained representations.

The use of RAID~\cite{dugan:2024} as an external training source addresses the domain-shift problem that has plagued previous PAN systems trained exclusively on in-task data. RAID's coverage of five adversarial attack types—none, homoglyph, synonym, paraphrase, and alternative spelling—across three content domains closely mirrors the obfuscation strategies anticipated in PAN 2026.
